\problem{1}
\textbf{[Отбор Сириус 2022, 8.5]}
У волшебника есть 5 одинаковых закрытых шкатулок, расставленных по кругу, в одной из которых лежит бриллиант. У Пети есть чашечные весы без гирь, которые честно укажут, что шкатулка с бриллиантом тяжелее, если она есть на весах. За один ход Петя может взять любые две шкатулки, взвесить их на чашечных весах. После взвешивания можно указать на любую шкатулку из пяти, если он уверен, что там есть бриллиант. Если же Петя не смог указать на шкатулку с бриллиантом, то волшебник произносит заклинание и бриллиант перемещается по часовой стрелке в соседнюю шкатулку. Если Петя ошибся, бриллиант исчезнет. Может ли Петя за два хода гарантированно найти шкатулку с бриллиантом?

\problem{2}
\textbf{[Отбор Сириус 2021, 8.1]}
Неотрицательные числа записаны в клетках квадрата 7 на 7. Сумма чисел в любых двух соседних столбцах не менее 20 , а в любых двух соседних строках не более 15 . Чему может быть равна сумма всех чисел, записанных в клетках данного квадрата?

\problem{3}
\textbf{[Отбор Сириус 2020, 8.3]}
Каждому из 10 человек выдали лист бумаги с напечатанными на нем цифрами от 0 до 9 и предложили вырезать несколько цифр. После этого оказалось, что если любые 9 человек объединят свои наборы вырезанных цифр, и все одинаковые цифры в наборе будут учитываться только один раз, то все полученные наборы окажутся различными и будут состоять из 9 цифр каждый (порядок цифр в наборе не учитывается). Можно ли определить, сколько каких цифр было вырезано?

\problem{4}
\textbf{[Отбор Сириус 2023, 8.6]}
В каждой ячейке таблицы $8 \times 8$ записано одно из чисел 1 или 2. Вычислили суммы чисел каждой строки и произведения чисел каждого столбца. Докажите, что из полученных 16-ти чисел есть два одинаковых.


\problem{5}
\textbf{[Отбор Сириус 2024, 8.5]}
На столе лежат десять тысяч спичек, а в корзине - две спички. Петя и Вася играют в игру. Они по очереди делают ходы, начинает Петя. Ход состоит в том, чтобы посчитать спички в корзине, пусть их будет $N$, выбрать делитель $d$ числа $N(d<N)$ и добавить в корзину $d$ спичек. Игра заканчивается, когда в корзине окажется более чем 2024 спички. Побеждает тот, кто сделал последний ход. У кого из игроков есть выигрышная стратегия?

\problem{6}
\textbf{[Отбор Сириус 2025, 8.6]}
Назовём две ладьи, размещенные на клетчатой доске «атакующими», если они находятся в одном столбце или в одной строке. Найдите наименьшее число $N$, такое, что если $N$ ладей размещены на доске $12 \times 12$, то всегда найдутся 7 ладей, таких, что никакие две не атакуют друг друга.